\subsection{Stationary innovation and covariance-matching controls}
\label{sec:innovation}
\InnovationStageNotice{}
At $\theta=0.10,0.25,0.40$, we compare shared and independent AR innovations
under the same weights and add covariance-compensated independent modes.
Both centred uniform and Gaussian innovations are used, giving six inputs
with 64 runs per cell. Every input is a circular length-$2T$ record with its
DC component removed. Here $T$ is the next 5-smooth integer at or above
$15L^{3/2}$, and the mode count is determined by $2T$.
Flat interfaces evolve under one record to $T$
and $2T$. Late-time squared widths are averaged over the last 20\% of regularly
stored times before taking the ensemble square root. This campaign controls
initialization and cutoff within its own comparisons; it does not replace
the cold-start protocol of the original shared-AR campaign.

For the positive weights in Eq.~\eqref{eq:ar}, define
\begin{align}
 A_n&=W_n\sum_m\frac{W_m}{1-r_nr_m},&
 V_n&=\sqrt{(1-r_n^2)A_n}.
 \label{eq:innovationmatch}
\end{align}
Independent innovations with weights $V_n$ reproduce the shared-input
stationary covariance at every lag, not just its variance.
The spectral identity survives sampling on a common frequency grid and
removing DC. Zero-mean Gaussian processes with the same complete covariance
have the same finite-dimensional laws; their thresholded fields and resulting
BD observables therefore also agree in law. The Gaussian comparison is an
implementation and sampling control, rather than evidence for an additional
physical effect of a latent mode representation.

Figure~\ref{fig:innovation}(a) and Table~\ref{tab:innovation} quantify how
innovation coupling and covariance compensation affect the effective width
slopes. Raw-covariance compensation cannot by itself match the actual
uniform-driven binary input: the maximum measured sign-correlation residual
across the control cells is
$\InnovationRawResidualMin$--$\InnovationRawResidualMax$.
Thresholding a non-Gaussian process depends on its two-time orthant
probabilities, which are not fixed by its raw covariance.

We therefore add a Gaussian surrogate fitted to the \emph{binary} covariance.
Only the first 32 uniform-shared noise records in each cell calibrate a
nonnegative mixture of circular AR spectra, using the Gaussian relation
$\rho(\tau)=\sin[\pi C_s(\tau)/2]$. Fresh Gaussian and uniform validation
records, with 64 runs and 16 sites per run, test every integer lag from 0
through 64 and the sampled long lags. A simultaneous bootstrap upper bound
on $\max_\tau|\Delta C_s(\tau)|$ must fall below 0.005, an explicit numerical
tolerance rather than a universal physical equivalence criterion.
Two initial cells, $\theta=0.40$ at $L=32,64$, failed this test.
We performed one documented calibration refinement, retaining the same
spectral family and tolerance: 128 fresh 64-site records calibrate every
integer lag from 0 through 64 and the original long lags, followed by
64 new, independent 64-site validation records per cell.
The initial failures motivated this refinement. After its calibration scheme
was fixed, the fresh validation records were not used for fitting, parameter
selection or further retries; the tolerance was not relaxed.
The simultaneous upper bounds changed as follows:
\InnovationRefinementBounds{}.
\InnovationLargeValidationNotice{}
With these refined controls, \InnovationBinaryPass{} of \InnovationCells{}
reported cells pass. The original failed archives remain in the numerical
supplement. \InnovationFailedInferenceNotice{}

The surrogate interface comparison uses the 32 held-out uniform-shared
runs. \InnovationSurrogateFinding{}
An unresolved difference does not prove equal exponents.
Independent four-time input statistics can nevertheless differ.
Sign-flip symmetry and stationarity imply
\begin{equation}
 \langle s_0s_1s_2s_3\rangle
 =16P(++++)-1-3C_s(1)-2C_s(2)-C_s(3).
 \label{eq:foursign}
\end{equation}
Figure~\ref{fig:innovation}(b) evaluates this moment with fresh validation
records at $L=256$, where all three covariance checks pass. The
$\theta=0.25$ and $0.40$ four-time input differences are resolved.
These observations distinguish a detectable high-order input change from a
demonstrated change in interface scaling. The original failed controls
are excluded from mechanistic inference. Only parameter grids passing
every retained covariance check are eligible for that comparison.
Unmeasured lag residuals and uncertainty in the fitted spectrum still
preclude a stronger attribution. Intervals use 2000 whole-run bootstrap draws,
preserve the six-input pairing within size, and resample sizes independently.
Surrogate intervals are conditional on the fitted spectrum; no
multiple-comparison adjustment is applied to these exploratory contrasts.
The raw run archives, original failures and reproducible analysis and
figure code are provided in the Supplemental Material~\cite{SupplementalData}.

\begin{table}[tb]
 \caption{Innovation-control slope differences at duration $2T$ across
 $L=\InnovationSizeRange$. G and U denote Gaussian and uniform innovations.
 Independent inputs are compared with their corresponding shared input;
 the binary-target Gaussian surrogate is compared with the held-out
 uniform-shared runs. Intervals are pointwise 95\% run-bootstrap intervals.}
 \label{tab:innovation}
 \compacttab
 \begin{ruledtabular}
 \begin{tabular}{clcc}
 $\theta$ & Input contrast & $\Delta\alpha_{\rm eff}$ & 95\% interval \\
 \InnovationContrastRows
 \end{tabular}
 \end{ruledtabular}
\end{table}

\begin{figure*}[tb]
 \centering
 \includegraphics[width=\textwidth]{fig7_innovation_controls}
 \caption{Stationary innovation controls.
 (a) Finite-size slope differences at $2T$ relative to the corresponding
 shared-input family. Each estimate uses \InnovationSizeCount{} sizes
 ($L=\InnovationSizeRange$) and 64 runs per input and size.
 (b) Gaussian binary-target minus fresh uniform-shared consecutive-four-sign
 moment at $L=256$, reconstructed by Eq.~\eqref{eq:foursign}; 64 independent
 validation runs per input, 16 sites per run. Bars are pointwise 95\%
 whole-run bootstrap intervals from 2000 draws.
 High-order input differences and finite-size slope sensitivity are distinct
 observables; approximate covariance matching does not establish identical
 full input processes or an asymptotic universality class.}
 \label{fig:innovation}
\end{figure*}
